%%%%%%%%%%%%%%%%%%%%%%%%%%%%%%%%%%%%%%%%%%%%%%%%%%%%%%%%%%%%%%%%%%%%%%
%% This template has been created from a previous Overleaf template
%% Author F. Brivio
%%%%%%%%%%%%%%%%%%%%%%%%%%%%%%%%%%%%%%%%%%%%%%%%%%%%%%%%%%%%%%%%%%%%%%
%% This is a (brief) model paper using the achemso class
%% The document class accepts keyval options, which should include
%% the target journal and optionally the manuscript type. 
%%%%%%%%%%%%%%%%%%%%%%%%%%%%%%%%%%%%%%%%%%%%%%%%%%%%%%%%%%%%%%%%%%%%%
%\documentclass[journal=jacsat,manuscript=article,email=false]{achemso}
\documentclass[10pt,a4paper]{article}
%\usepackage{achemso}
\usepackage[margin=1in]{geometry}
\usepackage[T1]{fontenc}
\usepackage[scaled]{helvet}
% Language setting
% Replace `english' with e.g. `spanish' to change the document language
\usepackage[english]{babel}

% Set page size and margins
% Replace `letterpaper' with `a4paper' for UK/EU standard size
%\usepackage[a4paper,top=2cm,bottom=2cm,left=3cm,right=3cm,marginparwidth=1.75cm]{geometry}

% Useful packages
\usepackage{amsmath}
\usepackage{graphicx}
\usepackage[colorlinks=true, allcolors=blue]{hyperref}
\usepackage[version=3]{mhchem}
\usepackage{setspace}
\title{Research Plan}
\author{Federico Brivio}
\date{}
\setlength{\voffset}{-0.75in}

%%%%%%%%%%%%%%%%%%%%%%%%%%%%%%%%%%%%%%%%%%%%%%%%%%%%%%%%%%%%%%%%%%%%%
%% Place any additional packages needed here.  Only include packages
%% which are essential, to avoid problems later. Do NOT use any
%% packages which require e-TeX (for example etoolbox): the e-TeX
%% extensions are not currently available on the ACS conversion
%% servers.
%%%%%%%%%%%%%%%%%%%%%%%%%%%%%%%%%%%%%%%%%%%%%%%%%%%%%%%%%%%%%%%%%%%%%
\usepackage[version=3]{mhchem} % Formula subscripts using \ce{}
%\setkeys{acs}{maxauthors=0}      % will list all authors
\usepackage{hyperref}
%\setkeys{acs}{articletitle=true}

%\newcommand*\mycommand[1]{\texttt{\emph{#1}}}
\usepackage[LGRgreek]{mathastext}    
\usepackage[
    sorting=none,
    giveninits=true,
    maxbibnames=100,
    style=numeric-comp,
    sortcites=true,
    backend=biber,
    defernumbers=true,
    natbib
]{biblatex}
%\addbibresource{achemso-demo.bib}
% Different bib files to be included
\addbibresource{list.bib}
\addbibresource{preprint.bib}
\addbibresource{thesis.bib}
% Add a specific  keyword to everything coming from a given .bib file
% In this example we used the keyword pre for the preprint.bif and thesis for thesis.bib
\DeclareSourcemap{
  \maps[datatype=bibtex]{
    \map[overwrite]{
      \perdatasource{preprint.bib}
      \step[fieldset=keywords, fieldvalue={,pre}, append]
    }
  }
}
\DeclareSourcemap{
  \maps[datatype=bibtex]{
    \map[overwrite]{
      \perdatasource{thesis.bib}
      \step[fieldset=keywords, fieldvalue={,thesis}, append]
    }
  }
}

\begin{document}
\section{Intruction - TO BE DELETED}
This document creates a list of all your publications divided per type and it is based on this page:
\href{https://tex.stackexchange.com/questions/358963/one-bibliography-section-per-bib-file-biblatex}{https://tex.stackexchange.com/questions/358963/one-bibliography-section-per-bib-file-biblatex} .

The publication are grouped by type in different files. In this case we distinguish between regular papers (list.bib), preprints ( preprint.bib) and the thesis (thesis.bib).

To list each type of publication independently, each file is marked with a specific tag.
This tag is \texttt{fieldvalue} and it is included in this code in the preamble:
\begin{verbatim}
\DeclareSourcemap{
  \maps[datatype=bibtex]{
    \map[overwrite]{
      \perdatasource{thesis.bib}
      \step[fieldset=keywords, fieldvalue={,thesis}, append]
    }
  }
}
\end{verbatim}
For each bib file it is possible to add a fieldvalue \texttt{fieldvalue={,thesis}} that allows to select which bib file is printed in a list and which is not.
To print a section that contains only the publications of a file with a specific tag the code is the following:
\begin{verbatim}
    \printbibliography[keyword={thesis},title=PhD Thesis]
\end{verbatim}
To print the papers not listed in a file with a specific tag:
\begin{verbatim}
    \printbibliography[notkeyword={pre},notkeyword={thesis},title=YOURNAME Publication list]
\end{verbatim}

Other divisions are of course possible. E.g. \textit{Publications as leading author}, \textit{Publications as last Author}, etc. .

A different method is available here: \href{https://bastian.rieck.me/blog/posts/2021/publication_lists/}{https://bastian.rieck.me/blog/posts/2021/publication\_lists}

\section{ACTUAL LISTS}
% For the final document cancell ALL from /begin{document) exclued to this line %%%%

%%%%%%%%%%%%%%%%%%%%%%%%%%%%%%%%%%%%%%%%%%%%%%%%%%%%%%%%%%%%%%%%%%%%%
%% The "tocentry" environment can be used to create an entry for the
%% graphical table of contents. It is given here as some journals
%% require that it is printed as part of the abstract page. It will
%% be automatically moved as appropriate.
%%%%%%%%%%%%%%%%%%%%%%%%%%%%%%%%%%%%%%%%%%%%%%%%%%%%%%%%%%%%%%%%%%%%%
\nocite{*}   % Since this is not a text document, we need to cite ALL the paper in the .bib file

\printbibliography[notkeyword={pre},notkeyword={thesis},title=YOURNAME Publication list]  % This will list all the publications execpt the ones contained in bib files tagged with the keywords pre and thesis

\printbibliography[keyword={pre},title=Preprint]  % This list only the publications contained in files tagged with the jeyword pre 
\printbibliography[keyword={thesis},title=PhD Thesis]
\end{document}